%% toc: P3 39 39.2
\section{Proofs: definitional complexity, rank definability and pebble games}
\label{sec:annex-p03-w03-p03-055-definitional-complexity}

Conventions are those of \cref{def:p03-w03-p03-055-setting}. The truth of a
formula of \(\mathrm{FO}\) or of \(\mathcal L_{\infty,\omega}\) under an
assignment depends only on the values of its free variables: this is an
induction on formulas in which the infinitary clauses are immediate, as
recorded at the start of \cref{sec:annex-p03-w03-p03-039-marker}. It is used
below without further mention.

\begin{annexitem}{annex:p03-w03-p03-055-rank-variables}{thm:p03-w03-p03-055-rank-variables}{Rank bounds the number of variables}
\begin{proof}
For a first-order formula \(\varphi\), a number \(d\ge0\) and an injective
map \(\rho\) from a finite set of variables containing the free variables
of \(\varphi\) into \(X_d\), define a formula
\(\operatorname{tr}(\varphi,\rho,d)\) by recursion on \(\varphi\):
\(\operatorname{tr}(\bottom,\rho,d)=\bottom\);
\(\operatorname{tr}(y=z,\rho,d)\) is \(\rho(y)=\rho(z)\);
\(\operatorname{tr}(R(y_1,\ldots,y_r),\rho,d)\) is
\(R(\rho(y_1),\ldots,\rho(y_r))\); \(\operatorname{tr}\) commutes with
\(\neg,\wedge,\vee\) (with the same \(\rho\) and \(d\), which is legitimate
because the free variables of a component are free in the compound); and
\[
 \operatorname{tr}(\exists y\,\psi,\rho,d)=\exists x_{d+1}\,
 \operatorname{tr}(\psi,\rho',d+1),\qquad
 \operatorname{tr}(\forall y\,\psi,\rho,d)=\forall x_{d+1}\,
 \operatorname{tr}(\psi,\rho',d+1),
\]
where \(\rho'\) agrees with \(\rho\) except that \(\rho'(y)=x_{d+1}\) (its
domain is \(\dom\rho\cup\{y\}\)). Since \(\rho\) maps into \(X_d\) and
\(x_{d+1}\notin X_d\), \(\rho'\) is injective into \(X_{d+1}\), and its
domain contains every free variable of \(\psi\).

By induction on \(\varphi\) we prove: \(\operatorname{tr}(\varphi,\rho,d)\)
has the quantifier rank of \(\varphi\), all its variables lie in
\(X_{d+\operatorname{qr}(\varphi)}\), and for every structure
\(\mathcal A\) and all assignments \(s,s'\) with \(s'(\rho(y))=s(y)\) for
\(y\in\dom\rho\),
\(\mathcal A\entails\varphi[s]\iff
\mathcal A\entails\operatorname{tr}(\varphi,\rho,d)[s']\).
For \(\bottom\) and atomic formulas the rank is \(0\), the variables are
values of \(\rho\), hence in \(X_d\), and the two sides say the same thing
about the same elements because \(s'(\rho(y))=s(y)\). The Boolean cases
follow directly from the induction hypothesis, since the rank of a compound
is the maximum of the ranks of its components. For \(\exists y\,\psi\):
\(\mathcal A\entails\exists y\,\psi[s]\) iff for some \(a\in A\),
\(\mathcal A\entails\psi[s[y\mapsto a]]\). The assignments
\(s[y\mapsto a]\) and \(s'[x_{d+1}\mapsto a]\) satisfy the hypothesis
for \(\rho'\): at \(y\) both give \(a\), and for \(z\in\dom\rho\setminus\{y\}\),
\(\rho'(z)=\rho(z)\in X_d\) differs from \(x_{d+1}\), so
\(s'[x_{d+1}\mapsto a](\rho'(z))=s'(\rho(z))=s(z)\). By the induction
hypothesis the condition is equivalent to
\(\mathcal A\entails\operatorname{tr}(\psi,\rho',d+1)[s'[x_{d+1}\mapsto a]]\)
for some \(a\), that is, to
\(\mathcal A\entails\operatorname{tr}(\exists y\,\psi,\rho,d)[s']\). The
rank rises by one on both sides, and the variables lie in
\(\{x_{d+1}\}\cup X_{d+1+\operatorname{qr}(\psi)}=
X_{d+\operatorname{qr}(\exists y\psi)}\). The case \(\forall y\,\psi\) is
the same computation with ``for every \(a\in A\)'' in place of ``for some
\(a\in A\)''.

For \(\varphi\) with free variables among \(x_1,\ldots,x_n\) take \(d=n\)
and \(\rho\) the identity on \(X_n\); then \(s'=s\) is allowed, so
\(\operatorname{tr}(\varphi,\rho,n)\) is equivalent to \(\varphi\), has the
same rank, and uses only variables in
\(X_{n+\operatorname{qr}(\varphi)}\subseteq X_{n+m}\). For a sentence take
\(n=0\) and \(\rho=\emptyset\): the variables lie in \(X_m\), so the
translation is an \(\mathrm{FO}^m\)-sentence.
\end{proof}
\end{annexitem}

\begin{annexitem}{annex:p03-w03-p03-055-rank-definability}{thm:p03-w03-p03-055-rank-definability}{Definability by rank and completeness of the game method}
\begin{proof}
(1) If \(\sigma\) has rank at most \(m\) and defines \(\mathcal K\) on
\(\mathcal C\), and \(\mathcal A,\mathcal B\in\mathcal C\) with
\(\mathcal A\equiv_m\mathcal B\) and \(\mathcal A\in\mathcal K\), then
\(\mathcal A\entails\sigma\), hence \(\mathcal B\entails\sigma\), hence
\(\mathcal B\in\mathcal K\).

Conversely assume the closure condition. Let
\(T=\{\operatorname{rec}_m(\mathcal A):\mathcal A\in\mathcal K\}\), the
records of length \(0\) and rank \(m\) of
\cref{def:p03-w03-p03-038-records}. It is a subset of the finite set
\(\mathrm{Rec}_{0,m}\) (\cref{thm:p03-w03-p03-038-finite-records}), so
\(\sigma\defeq\bigvee_{t\in T}\chi_t\) is a finite disjunction of the
characteristic sentences of \cref{def:p03-w03-p03-038-characteristic}; it
is \(\bottom\) if \(T=\emptyset\). Each \(\chi_t\) is a sentence of rank at
most \(m\) (\cref{thm:p03-w03-p03-038-characteristic}), so \(\sigma\) is a
first-order sentence of rank at most \(m\). Let
\(\mathcal B\in\mathcal C\). If \(\mathcal B\in\mathcal K\), then
\(t=\operatorname{rec}_m(\mathcal B)\in T\) and
\(\mathcal B\entails\chi_t\) by \cref{thm:p03-w03-p03-038-characteristic},
so \(\mathcal B\entails\sigma\). If \(\mathcal B\entails\sigma\), then
\(\mathcal B\entails\chi_t\) for some \(t=\operatorname{rec}_m(\mathcal A)\)
with \(\mathcal A\in\mathcal K\); by
\cref{thm:p03-w03-p03-038-characteristic},
\(\operatorname{rec}_m(\mathcal B)=t=\operatorname{rec}_m(\mathcal A)\), and
by \cref{thm:p03-w03-p03-038-pointed-ef} (with \(n=0\), implication from
equal records to agreement) \(\mathcal A\equiv_m\mathcal B\). The closure
condition gives \(\mathcal B\in\mathcal K\).

(2) Every first-order sentence has some finite rank, so by (1)
\(\mathcal K\) is first-order definable on \(\mathcal C\) iff the closure
condition of (1) holds for some \(m\). The condition fails for \(m\)
exactly when there are \(\mathcal A\in\mathcal K\) and
\(\mathcal B\in\mathcal C\setminus\mathcal K\) with
\(\mathcal A\equiv_m\mathcal B\), and by
\cref{thm:p03-w03-p03-038-pointed-ef} (with \(n=0\)) this is the same as
Duplicator having a winning strategy in the \(m\)-round game on them. Hence
\(\mathcal K\) is not first-order definable on \(\mathcal C\) iff such a
pair exists for every \(m\).
\end{proof}
\end{annexitem}

\begin{annexitem}{annex:p03-w03-p03-055-pebble}{thm:p03-w03-p03-055-pebble}{Pebble games characterize \(k\)-variable equivalence}
\begin{proof}
\emph{(1)\(\Rightarrow\)(2).} Let \(\mathcal I\) be a winning family with
\(\langle s,t\rangle\in\mathcal I\). We prove by induction on
\(\mathcal L^k_{\infty\omega}\)-formulas \(\varphi\): for all
\(s',t'\) with a common domain \(D'\subseteq X_k\) containing the free
variables of \(\varphi\) and with \(\langle s',t'\rangle\in\mathcal I\),
\(\mathcal A\entails\varphi[s']\iff\mathcal B\entails\varphi[t']\).
Both sides of \(\bottom\) fail. For \(y=z\) with \(y,z\in D'\), the
pairs \((s'(y),t'(y))\) and \((s'(z),t'(z))\) belong to the partial
isomorphism \(\langle s',t'\rangle\), so \(s'(y)=s'(z)\iff t'(y)=t'(z)\);
for \(R(y_1,\ldots,y_r)\) the pairs \((s'(y_i),t'(y_i))\) belong to it, so
\(R\) holds of the \(s'\)-values iff it holds of the \(t'\)-values. For \(\neg\varphi\), \(\bigwedge\Phi\) and \(\bigvee\Phi\) the
free variables of each component lie in \(D'\), and the induction
hypothesis applied to every component gives the claim.

Let \(\varphi=\exists x_i\,\psi\), put \(D''=D'\setminus\{x_i\}\) and
\(f=\langle s'|_{D''},t'|_{D''}\rangle\). Then
\(f\subseteq\langle s',t'\rangle\), so \(f\in\mathcal I\), and
\(|f|\le|D''|\le k-1<k\) because \(D''\subseteq X_k\setminus\{x_i\}\).
Suppose \(\mathcal A\entails\varphi[s']\) with witness \(a\), so
\(\mathcal A\entails\psi[s'|_{D''}[x_i\mapsto a]]\) (this assignment is
\(s'[x_i\mapsto a]\)). By the forth property there is \(g\in\mathcal I\)
with \(f\subseteq g\) and \(a\in\dom g\); let \(b=g(a)\). Then
\(\langle s'|_{D''}[x_i\mapsto a],t'|_{D''}[x_i\mapsto b]\rangle=
f\cup\{(a,b)\}\subseteq g\) lies in \(\mathcal I\), and the free variables
of \(\psi\) lie in \(D''\cup\{x_i\}\). The induction hypothesis gives
\(\mathcal B\entails\psi[t'|_{D''}[x_i\mapsto b]]\), so
\(\mathcal B\entails\varphi[t']\). The implication from \(\mathcal B\) to
\(\mathcal A\) uses the back property to find \(a\) for a witness \(b\).
For \(\varphi=\forall x_i\,\psi\) with \(\mathcal A\entails\varphi[s']\):
for any \(b\in B\) the back property gives \(g\supseteq f\) in
\(\mathcal I\) with \(b\) in its range; with \(a=g^{-1}(b)\), the relation
\(f\cup\{(a,b)\}\subseteq g\) is in \(\mathcal I\), and
\(\mathcal A\entails\psi[s'[x_i\mapsto a]]\) gives
\(\mathcal B\entails\psi[t'[x_i\mapsto b]]\); so
\(\mathcal B\entails\varphi[t']\). The converse uses the forth property.
Applied to \(s,t\), the claim gives
\((\mathcal A,s)\equiv^k_{\infty\omega}(\mathcal B,t)\).

\emph{(2)\(\Rightarrow\)(1).} Let \(\mathcal I\) be the set of all relations
\(\langle s',t'\rangle\) with \(s'\colon D'\to A\), \(t'\colon D'\to B\),
\(D'\subseteq X_k\) and \((\mathcal A,s')\equiv^k_{\infty\omega}(\mathcal B,t')\).
It contains \(\langle s,t\rangle\). Each member is a partial
isomorphism: the formulas \(y=z\) and \(R(y_1,\ldots,y_r)\) with variables
in \(D'\) belong to \(\mathcal L^k_{\infty\omega}\), and agreement on them
is exactly the defining condition. Each member has at most
\(|D'|\le k\) pairs.

Closure under subsets: let \(f\subseteq\langle s',t'\rangle\) and
\(D''=\{x\in D':(s'(x),t'(x))\in f\}\). Every pair of \(f\) has the form
\((s'(x),t'(x))\) with \(x\in D''\), so
\(f=\langle s'|_{D''},t'|_{D''}\rangle\); formulas with free variables in
\(D''\) have them in \(D'\), so
\((\mathcal A,s'|_{D''})\equiv^k_{\infty\omega}(\mathcal B,t'|_{D''})\) and
\(f\in\mathcal I\).

Forth: let \(f=\langle s',t'\rangle\in\mathcal I\) with \(|f|<k\), and
\(a\in A\). Choose \(D''\subseteq D'\) containing, for each pair in \(f\),
exactly one variable \(x\) with \((s'(x),t'(x))\) equal to that pair. Then
\(f=\langle s'',t''\rangle\) for \(s''=s'|_{D''}\), \(t''=t'|_{D''}\), the
two restrictions are still \(\equiv^k_{\infty\omega}\)-equivalent, and
\(|D''|=|f|<k\); choose \(x_i\in X_k\setminus D''\). We claim that some
\(b\in B\) satisfies
\((\mathcal A,s''[x_i\mapsto a])\equiv^k_{\infty\omega}(\mathcal B,t''[x_i\mapsto b])\).
Otherwise choose for each \(b\in B\) an
\(\mathcal L^k_{\infty\omega}\)-formula \(\theta_b\) with free variables in
\(D''\cup\{x_i\}\) such that \(\mathcal A\entails\theta_b[s''[x_i\mapsto a]]\)
and \(\mathcal B\not\entails\theta_b[t''[x_i\mapsto b]]\) (replacing a
formula by its negation if the disagreement goes the other way). The formula
\(\Theta=\exists x_i\bigwedge_{b\in B}\theta_b\) uses only variables of
\(X_k\), so it lies in \(\mathcal L^k_{\infty\omega}\), and its free
variables lie in \(D''\). It holds in \(\mathcal A\) under \(s''\), with
witness \(a\), and fails in \(\mathcal B\) under \(t''\), since every
\(b\) falsifies \(\theta_b\). This contradicts the equivalence of
\((\mathcal A,s'')\) and \((\mathcal B,t'')\). With such a \(b\),
\(g=\langle s''[x_i\mapsto a],t''[x_i\mapsto b]\rangle=f\cup\{(a,b)\}\) is
in \(\mathcal I\), extends \(f\) and has \(a\) in its domain.

Back: with \(f\), \(D''\), \(x_i\) as before and \(b\in B\), suppose no
\(a\in A\) gives
\((\mathcal A,s''[x_i\mapsto a])\equiv^k_{\infty\omega}(\mathcal B,t''[x_i\mapsto b])\).
Choose for each \(a\in A\) a formula \(\theta_a\) of
\(\mathcal L^k_{\infty\omega}\) with free variables in \(D''\cup\{x_i\}\),
true in \(\mathcal B\) under \(t''[x_i\mapsto b]\) and false in
\(\mathcal A\) under \(s''[x_i\mapsto a]\). Then
\(\exists x_i\bigwedge_{a\in A}\theta_a\) holds in \(\mathcal B\) under
\(t''\) and fails in \(\mathcal A\) under \(s''\), a contradiction; the
\(a\) so obtained gives \(g=f\cup\{(a,b)\}\in\mathcal I\) with \(b\) in its
range. So \(\mathcal I\) is a winning family containing
\(\langle s,t\rangle\).

\emph{Finite structures.} (2)\(\Rightarrow\)(3) holds because
\(\mathrm{FO}^k\subseteq\mathcal L^k_{\infty\omega}\). For
(3)\(\Rightarrow\)(1), let \(\mathcal I_{\mathrm{FO}}\) be defined as
\(\mathcal I\) above with \(\equiv^k\) in place of
\(\equiv^k_{\infty\omega}\). The verification that its members are partial
isomorphisms with at most \(k\) pairs, and that it is closed under subsets,
used only atomic formulas and the dependence of truth on free variables, and
applies verbatim to \(\mathrm{FO}^k\). In the forth step the formulas
\(\theta_b\) are now chosen in \(\mathrm{FO}^k\); since \(B\) is finite,
\(\bigwedge_{b\in B}\theta_b\) is a finite conjunction, so
\(\Theta\in\mathrm{FO}^k\), and the contradiction is with \(\equiv^k\). In
the back step the conjunction over the finite set \(A\) is finite for the
same reason. Hence \(\mathcal I_{\mathrm{FO}}\) is a winning family
containing \(\langle s,t\rangle\). Together with (1)\(\Rightarrow\)(2) this
closes the cycle of implications.
\end{proof}
\end{annexitem}

\begin{annexitem}{annex:p03-w03-p03-055-k-definability}{thm:p03-w03-p03-055-k-definability}{Definability in \(k\)-variable infinitary logic on finite structures}
\begin{proof}
For finite structures, \(\mathcal A\equiv^k\mathcal B\) is equivalent to
Duplicator winning the \(k\)-pebble game and to
\(\mathcal A\equiv^k_{\infty\omega}\mathcal B\), by
\cref{thm:p03-w03-p03-055-pebble} with \(D=\emptyset\).

Suppose an \(\mathcal L^k_{\infty\omega}\)-sentence \(\Phi\) defines
\(\mathcal K\) on \(\mathcal C\). If \(\mathcal A,\mathcal B\in\mathcal C\),
\(\mathcal A\equiv^k\mathcal B\) and \(\mathcal A\in\mathcal K\), then
\(\mathcal A\equiv^k_{\infty\omega}\mathcal B\), so \(\mathcal B\entails\Phi\)
and \(\mathcal B\in\mathcal K\).

Conversely assume the closure condition. For a finite structure
\(\mathcal A\) let \(\sigma_{\mathcal A}\) be the conjunction of the set of
all \(\mathrm{FO}^k\)-sentences true in \(\mathcal A\); it is an
\(\mathcal L^k_{\infty\omega}\)-sentence. A structure satisfies
\(\sigma_{\mathcal A}\) iff it is \(\equiv^k\)-equivalent to \(\mathcal A\):
if \(\mathcal B\entails\sigma_{\mathcal A}\) and \(\theta\) is an
\(\mathrm{FO}^k\)-sentence, then \(\theta\) or \(\neg\theta\) is a conjunct,
according to its truth in \(\mathcal A\), and \(\mathcal B\) agrees; the
converse is immediate. Let \(R\) be the set of members of \(\mathcal K\)
whose universe is \(\{0,\ldots,n-1\}\) for some \(n\ge1\); it is a set, and
every member of \(\mathcal K\) is isomorphic to one of its members. Put
\(\Phi=\bigvee_{\mathcal A\in R}\sigma_{\mathcal A}\). For
\(\mathcal B\in\mathcal C\): if \(\mathcal B\in\mathcal K\), choose
\(\mathcal A\in R\) isomorphic to \(\mathcal B\); isomorphic structures
satisfy the same sentences (\cref{thm:p03-w03-p03-039-iso-invariance}), so
\(\mathcal B\entails\sigma_{\mathcal A}\) and \(\mathcal B\entails\Phi\). If
\(\mathcal B\entails\Phi\), then \(\mathcal B\equiv^k\mathcal A\) for some
\(\mathcal A\in R\subseteq\mathcal K\), and the closure condition gives
\(\mathcal B\in\mathcal K\).

Finally \(\mathcal L^\omega_{\infty\omega}\) is the union of the
\(\mathcal L^k_{\infty\omega}\), so a defining sentence of it lies in some
\(\mathcal L^k_{\infty\omega}\), and the first part applies with that \(k\).
\end{proof}
\end{annexitem}

\begin{annexitem}{annex:p03-w03-p03-055-unbounded}{thm:p03-w03-p03-055-unbounded}{Without a variable bound every finite-structure class is definable}
\begin{proof}
Let \(\mathcal A\) be finite with universe \(\{a_1,\ldots,a_n\}\), the
\(a_i\) distinct, \(\bar a=(a_1,\ldots,a_n)\), and
\(S=\operatorname{rec}_0(\mathcal A,\bar a)\subseteq\mathrm{At}_n\), with
\(\chi_S\) as in \cref{def:p03-w03-p03-038-characteristic}. Put
\[
 \delta_{\mathcal A}\defeq\exists x_1\cdots\exists x_n\Bigl(
 \chi_S\wedge
 \forall x_{n+1}\bigvee_{i=1}^n x_{n+1}=x_i\Bigr).
\]
If \(h\colon\mathcal A\to\mathcal B\) is an isomorphism, the elements
\(h(a_i)\) witness \(\delta_{\mathcal A}\) in \(\mathcal B\): an isomorphism
preserves the truth of atomic formulas, so the \(h(a_i)\) satisfy exactly
the members of \(S\) among the formulas of \(\mathrm{At}_n\), and they
exhaust \(B\). Conversely, if \(b_1,\ldots,b_n\) witness \(\delta_{\mathcal A}\) in
\(\mathcal B\), then \(\langle\bar a,\bar b\rangle\) is a partial
isomorphism: by \cref{thm:p03-w03-p03-038-characteristic}
with \(j=0\), \(\operatorname{rec}_0(\mathcal B,\bar b)=S\), so the two
tuples satisfy the same formulas of \(\mathrm{At}_n\), which is the
definition of a partial isomorphism in
\cref{def:p03-w03-p03-038-positions}; whose domain is all of \(A\) and whose range is all of \(B\) by the
last conjunct; so it is an isomorphism. Thus \(\mathcal B\entails
\delta_{\mathcal A}\) iff \(\mathcal B\cong\mathcal A\).

Since \(\tau\) is finite, for each \(n\) there are finitely many
\(\tau\)-structures with universe \(\{0,\ldots,n-1\}\), so the set \(R\) of
members of \(\mathcal K\) with such a universe is countable, and every
member of \(\mathcal K\) is isomorphic to a member of \(R\). The countable
disjunction \(\bigvee_{\mathcal A\in R}\delta_{\mathcal A}\) of first-order
sentences holds in a finite \(\mathcal B\) iff \(\mathcal B\) is isomorphic
to a member of \(R\), iff \(\mathcal B\in\mathcal K\), because
\(\mathcal K\) is closed under isomorphism.
\end{proof}
\end{annexitem}

\begin{annexitem}{annex:p03-w03-p03-055-pebble-basic}{thm:p03-w03-p03-055-pebble-basic}{Monotonicity in the number of pebbles}
\begin{proof}
Spoiler wins the \(k\)-pebble game exactly when Duplicator has no winning
family in the sense of \cref{def:p03-w03-p03-055-pebble}, so the two
formulations of the first claim are contrapositives of each other. By
\cref{thm:p03-w03-p03-055-pebble} with \(D=\emptyset\), Duplicator wins the
\(k\)-pebble game on \(\mathcal A\) and \(\mathcal B\) if and only if
\(\mathcal A\equiv^k_{\infty\omega}\mathcal B\), and likewise for \(k'\).
Since \(X_k\subseteq X_{k'}\), every \(\mathcal L^k_{\infty\omega}\)-sentence
is an \(\mathcal L^{k'}_{\infty\omega}\)-sentence. Hence
\(\mathcal A\equiv^{k'}_{\infty\omega}\mathcal B\) implies
\(\mathcal A\equiv^k_{\infty\omega}\mathcal B\), which is the first claim.
The relation \(\equiv^k_{\infty\omega}\) between structures says that every
\(\mathcal L^k_{\infty\omega}\)-sentence has the same truth value in both. Equality of truth values is
reflexive, symmetric and transitive, sentence by sentence, so
\(\equiv^k_{\infty\omega}\) is an equivalence relation, and by
\cref{thm:p03-w03-p03-055-pebble} so is the relation of the game.
\end{proof}
\end{annexitem}

\begin{annexitem}{annex:p03-w03-p03-055-finitely-many-k}{thm:p03-w03-p03-055-finitely-many-k}{Finitely many \(k\)-variable formulas of each rank}
\begin{proof}
For \(E\subseteq X_k\) let \(F_m(E)\) be the set of
\(\mathrm{FO}^k\)-formulas of quantifier rank at most \(m\) whose free
variables lie in \(E\). For formulas with free variables in \(X_k\) write
\(\varphi\sim\psi\) when \(\mathcal A\entails\varphi[s]\iff
\mathcal A\entails\psi[s]\) for every structure \(\mathcal A\) and every
\(s\colon X_k\to A\). This is logical equivalence: an assignment whose
domain contains the free variables of both formulas can be restricted to
those variables and then extended to all of \(X_k\), as \(A\neq\emptyset\),
without changing either truth value.

Call \emph{basic of level \(0\)} the atomic formulas with variables in
\(X_k\), and \emph{basic of level \(m+1\)} these together with the formulas
\(\exists x_i\,\psi\) and \(\forall x_i\,\psi\) for \(i\le k\) and
\(\psi\in F_m(X_k)\). By induction on formulas, every member of
\(F_0(X_k)\) is built from basic formulas of level \(0\) and \(\bottom\) by
\(\neg,\wedge,\vee\) (a formula of rank \(0\) has no quantifier), and every
member of \(F_{m+1}(X_k)\) is built in the same way from basic formulas of
level \(m+1\) and \(\bottom\): a quantified subformula
\(\exists x_i\,\psi\) or \(\forall x_i\,\psi\) met while descending through
the connectives has rank at most \(m+1\) and variables in \(X_k\), so
\(\psi\in F_m(X_k)\).

We show by induction on \(m\) that \(F_m(X_k)\) meets only finitely many
\(\sim\)-classes. There are \(k^2\) equalities and, for each
\(R\in\tau\) of arity \(r\), \(k^r\) atomic formulas \(R(y_1,\ldots,y_r)\)
with \(y_j\in X_k\); as \(\tau\) is finite there are finitely many atomic
formulas with variables in \(X_k\). For level \(m+1\), let
\(\psi_1,\ldots,\psi_N\) represent the classes met by \(F_m(X_k)\) (the
induction hypothesis). If \(\psi\sim\psi_j\), then
\(\exists x_i\,\psi\sim\exists x_i\,\psi_j\): for \(s\colon X_k\to A\),
\(\mathcal A\entails\exists x_i\,\psi[s]\) iff
\(\mathcal A\entails\psi[s[x_i\mapsto a]]\) for some \(a\in A\), and
\(s[x_i\mapsto a]\) is again an assignment on \(X_k\), so this holds iff
\(\mathcal A\entails\psi_j[s[x_i\mapsto a]]\) for some \(a\), iff
\(\mathcal A\entails\exists x_i\,\psi_j[s]\); the universal case is the same
with ``for every \(a\)''. So the basic formulas of each level fall into
finitely many classes; let \(\beta_1,\ldots,\beta_M\) represent them.

For \(\mathcal A\) and \(s\colon X_k\to A\) let
\(v(\mathcal A,s)\in\{0,1\}^M\) list the truth values of
\(\beta_1,\ldots,\beta_M\). To a formula \(\varphi\) built from basic
formulas and \(\bottom\) by connectives assign a function
\(f_\varphi\colon\{0,1\}^M\to\{0,1\}\): for a basic formula, the projection
onto the coordinate of its representative; \(f_\bottom=0\); and \(f\)
commutes with the connectives, read as operations on \(\{0,1\}\). By
induction on this construction, the truth value of \(\varphi\) at
\((\mathcal A,s)\) is \(f_\varphi(v(\mathcal A,s))\). Hence formulas with the
same function are \(\sim\)-equivalent, and the classes met by the level are
at most the \(2^{2^M}\) functions. This completes the induction.

Finally, \(F_m(D)\subseteq F_m(X_k)\), so \(F_m(D)\) meets only finitely
many classes; choosing one member of \(F_m(D)\) in each class it meets gives
a finite set of representatives.
\end{proof}
\end{annexitem}

\begin{annexitem}{annex:p03-w03-p03-055-bounded-pebble}{thm:p03-w03-p03-055-bounded-pebble}{The \(m\)-round \(k\)-pebble game characterizes \(\equiv^k_m\)}
\begin{proof}
\emph{Game implies agreement.} We prove by induction on the
\(\mathrm{FO}^k\)-formula \(\varphi\): for every \(m\ge\operatorname{qr}(\varphi)\)
and every position \((s',t')\) whose common domain \(D'\) contains the free
variables of \(\varphi\) and from which Duplicator wins
\(G^k_m(\mathcal A,s';\mathcal B,t')\),
\(\mathcal A\entails\varphi[s']\iff\mathcal B\entails\varphi[t']\). In
either clause of \cref{def:p03-w03-p03-055-bounded-pebble} a win requires
\(\langle s',t'\rangle\) to be a partial isomorphism. Both sides of
\(\bottom\) fail. For \(y=z\) with \(y,z\in D'\) the pairs
\((s'(y),t'(y))\) and \((s'(z),t'(z))\) lie in \(\langle s',t'\rangle\), so
\(s'(y)=s'(z)\iff t'(y)=t'(z)\); for \(R(y_1,\ldots,y_r)\) the pairs
\((s'(y_j),t'(y_j))\) lie in it, so \(R\) holds of the \(s'\)-values iff it
holds of the \(t'\)-values. For \(\neg,\wedge,\vee\) the components have rank
at most \(m\) and free variables in \(D'\), and the induction hypothesis
applies with the same \(m\) and position.

Let \(\varphi=\exists x_i\,\psi\), so \(m\ge1\); write \(m=m'+1\), with
\(m'\ge\operatorname{qr}(\psi)\). Suppose \(\mathcal A\entails\varphi[s']\),
witnessed by \(a\): \(\mathcal A\entails\psi[s'[x_i\mapsto a]]\). Clause 2
of the definition, applied to the label \(i\) and the element \(a\), gives
\(b\in B\) such that Duplicator wins
\(G^k_{m'}(\mathcal A,s'[x_i\mapsto a];\mathcal B,t'[x_i\mapsto b])\). The new
common domain \(D'\cup\{x_i\}\) contains the free variables of \(\psi\), so
the induction hypothesis gives \(\mathcal B\entails\psi[t'[x_i\mapsto b]]\)
and \(\mathcal B\entails\varphi[t']\). The converse implication uses the
clause for moves in \(\mathcal B\). For \(\varphi=\forall x_i\,\psi\) and
\(\mathcal A\entails\varphi[s']\), let \(b\in B\); the clause for moves in
\(\mathcal B\) supplies \(a\) with a winning position
\((s'[x_i\mapsto a],t'[x_i\mapsto b])\) for \(m'\) rounds; since
\(\psi\) holds at \(s'[x_i\mapsto a]\), the induction hypothesis makes it hold
at \(t'[x_i\mapsto b]\). So \(\mathcal B\entails\varphi[t']\), and the
converse uses the clause for moves in \(\mathcal A\). Applied to \((s,t)\)
and the formulas of rank at most \(m\) with free variables in \(D\), this
gives \((\mathcal A,s)\equiv^k_m(\mathcal B,t)\).

\emph{Agreement implies game.} By induction on \(m\), for all positions.
Agreement on the atomic formulas with variables in \(D\) makes
\(\langle s,t\rangle\) a partial isomorphism: for \(x,y\in D\) the formula
\(x=y\) gives \(s(x)=s(y)\iff t(x)=t(y)\), and the formulas
\(R(y_1,\ldots,y_r)\) with \(y_j\in D\) give preservation and reflection of
\(R\). Atomic formulas have rank \(0\le m\). This settles \(m=0\) and the
first requirement of clause 2.

Let \((\mathcal A,s)\equiv^k_{m+1}(\mathcal B,t)\), let \(i\le k\) and
\(a\in A\), and put \(D'=D\cup\{x_i\}\). Suppose that for no \(b\in B\) is
\((\mathcal A,s[x_i\mapsto a])\equiv^k_m(\mathcal B,t[x_i\mapsto b])\). By
\cref{thm:p03-w03-p03-055-finitely-many-k} fix a finite set \(\Psi\) of
\(\mathrm{FO}^k\)-formulas of rank at most \(m\) with free variables in
\(D'\) representing every such formula up to logical equivalence; \(\Psi\)
is closed under negation up to equivalence, since a negation has the same
rank, variables and free variables. For each \(b\in B\) there is a formula
of rank at most \(m\), with free variables in \(D'\), true at
\(s[x_i\mapsto a]\) and false at \(t[x_i\mapsto b]\) (if the disagreement goes
the other way, take the negation); let \(\theta_b\in\Psi\) be equivalent to
it. Then \(\Theta\defeq\exists x_i\bigwedge\{\theta_b:b\in B\}\) is a finite
conjunction under one quantifier, an \(\mathrm{FO}^k\)-formula of rank at most
\(m+1\) whose free variables lie in \(D'\setminus\{x_i\}\subseteq D\). It
holds in \(\mathcal A\) at \(s\), with witness \(a\), and fails in
\(\mathcal B\) at \(t\), since every \(b\) falsifies \(\theta_b\) at
\(t[x_i\mapsto b]\). This contradicts
\((\mathcal A,s)\equiv^k_{m+1}(\mathcal B,t)\). So some \(b\) gives
\((\mathcal A,s[x_i\mapsto a])\equiv^k_m(\mathcal B,t[x_i\mapsto b])\), and by
the induction hypothesis Duplicator wins \(G^k_m\) from that position. For
\(i\le k\) and \(b\in B\) the same argument with the roles of the
structures exchanged, using \(\exists x_i\bigwedge\{\theta_a:a\in A\}\),
gives \(a\). All requirements of clause 2 hold, and Duplicator wins
\(G^k_{m+1}(\mathcal A,s;\mathcal B,t)\).
\end{proof}
\end{annexitem}

\begin{annexitem}{annex:p03-w03-p03-055-bounded-consequences}{thm:p03-w03-p03-055-bounded-consequences}{Consequences of the two-parameter game}
\begin{proof}
(1) Every \(\mathrm{FO}^k\)-formula has some finite quantifier rank, so
\((\mathcal A,s)\equiv^k(\mathcal B,t)\) holds if and only if
\((\mathcal A,s)\equiv^k_m(\mathcal B,t)\) for every \(m\), and by
\cref{thm:p03-w03-p03-055-bounded-pebble} if and only if Duplicator wins
every \(G^k_m(\mathcal A,s;\mathcal B,t)\). For finite structures,
\cref{thm:p03-w03-p03-055-pebble} identifies \(\equiv^k\) with Duplicator
winning the \(k\)-pebble game.

(2) Every \(\mathrm{FO}^k\)-sentence of rank at most \(m\) is a first-order
sentence of rank at most \(m\), so \(\mathcal A\equiv_m\mathcal B\) implies
\(\mathcal A\equiv^k_m\mathcal B\). Conversely let \(\sigma\) be a
first-order sentence of rank at most \(m\). If \(m\ge1\),
\cref{thm:p03-w03-p03-055-rank-variables} gives an equivalent sentence of the
same rank with variables in \(X_m\subseteq X_k\), on which
\(\mathcal A\) and \(\mathcal B\) agree when
\(\mathcal A\equiv^k_m\mathcal B\); so they agree on \(\sigma\). If
\(m=0\), \(\sigma\) has no quantifier and no free variable, so it has no
variable at all; every atomic formula contains a variable, as all symbols
of \(\tau\) have positive arity, so \(\sigma\) is built from \(\bottom\) by
connectives and is itself an \(\mathrm{FO}^k\)-sentence of rank \(0\). Hence
\(\mathcal A\equiv^k_m\mathcal B\) implies \(\mathcal A\equiv_m\mathcal B\).
The equivalence with the \(m\)-round Ehrenfeucht--Fra\"iss\'e game is
\cref{thm:p03-w03-p03-038-pointed-ef} for tuples of length \(0\).

(3) If an \(\mathrm{FO}^k\)-sentence \(\sigma\) of rank at most \(m\) defines
\(\mathcal K\) on \(\mathcal C\), \(\mathcal A\in\mathcal K\) and
\(\mathcal A\equiv^k_m\mathcal B\) with \(\mathcal B\in\mathcal C\), then
\(\mathcal B\entails\sigma\) and \(\mathcal B\in\mathcal K\). Conversely
assume the closure condition. By
\cref{thm:p03-w03-p03-055-finitely-many-k} with \(D=\emptyset\) there are
\(\mathrm{FO}^k\)-sentences \(\theta_1,\ldots,\theta_N\) of rank at most
\(m\) such that every \(\mathrm{FO}^k\)-sentence of rank at most \(m\) is
logically equivalent to one of them. For a structure \(\mathcal A\) let
\[
 \zeta_{\mathcal A}\defeq\bigwedge_{\mathcal A\entails\theta_j}\theta_j\wedge
 \bigwedge_{\mathcal A\not\entails\theta_j}\neg\theta_j ,
\]
an \(\mathrm{FO}^k\)-sentence of rank at most \(m\). A structure
\(\mathcal B\) satisfies \(\zeta_{\mathcal A}\) iff
\(\mathcal B\equiv^k_m\mathcal A\): if \(\mathcal B\entails\zeta_{\mathcal A}\),
then \(\mathcal B\) agrees with \(\mathcal A\) on every \(\theta_j\), hence on
every \(\mathrm{FO}^k\)-sentence of rank at most \(m\), each being equivalent
to some \(\theta_j\); the converse is immediate. The sentence
\(\zeta_{\mathcal A}\) is determined by the set of \(j\) with
\(\mathcal A\entails\theta_j\), so as \(\mathcal A\) ranges over
\(\mathcal K\) at most \(2^N\) distinct sentences occur; let \(\sigma\) be
their disjunction (\(\bottom\) if \(\mathcal K\) is empty), an
\(\mathrm{FO}^k\)-sentence of rank at most \(m\). If
\(\mathcal B\in\mathcal K\) then \(\mathcal B\entails\zeta_{\mathcal B}\), a
disjunct. If \(\mathcal B\in\mathcal C\) satisfies \(\sigma\), then
\(\mathcal B\entails\zeta_{\mathcal A}\) for some \(\mathcal A\in\mathcal K\), so
\(\mathcal A\equiv^k_m\mathcal B\) and the closure condition gives
\(\mathcal B\in\mathcal K\).

Every \(\mathrm{FO}^k\)-sentence has a finite rank, so \(\mathcal K\) is
\(\mathrm{FO}^k\)-definable on \(\mathcal C\) iff the closure condition holds
for some \(m\). It fails for \(m\) exactly when some
\(\mathcal A_m\in\mathcal K\) and \(\mathcal B_m\in\mathcal C\setminus
\mathcal K\) satisfy \(\mathcal A_m\equiv^k_m\mathcal B_m\), that is, by
\cref{thm:p03-w03-p03-055-bounded-pebble} with \(D=\emptyset\), when
Duplicator wins \(G^k_m(\mathcal A_m,\emptyset;\mathcal B_m,\emptyset)\).
\end{proof}
\end{annexitem}

\begin{annexitem}{annex:p03-w03-p03-055-cliques}{thm:p03-w03-p03-055-cliques}{Cliques and the strictness of both hierarchies}
\begin{proof}
(1) In a clique two vertices are adjacent iff they are distinct. Hence an
injective partial map between cliques preserves and reflects adjacency and
equality, and is a partial isomorphism. Let \(\mathcal I\) be the set of
injective partial maps \(f\) from \(K_n\) to \(K_{n'}\) with \(|f|\le k\).
It contains \(\emptyset\) and is closed under subsets. Forth: let
\(f\in\mathcal I\), \(|f|<k\), and let \(a\) be a vertex of \(K_n\). If
\(a\in\dom f\), take \(g=f\). Otherwise the range of \(f\) has fewer than
\(k\le n'\) elements, so some vertex \(b\) of \(K_{n'}\) is outside it, and
\(g=f\cup\{(a,b)\}\) is injective with \(|g|\le k\). Back: for a vertex
\(b\) of \(K_{n'}\) not in the range of \(f\), the domain of \(f\) has fewer
than \(k\le n\) elements, and any vertex \(a\) of \(K_n\) outside it gives
\(f\cup\{(a,b)\}\in\mathcal I\); if \(b\) is in the range, take \(g=f\).
So \(\mathcal I\) is a winning family and
\cref{thm:p03-w03-p03-055-pebble} (with \(D=\emptyset\)) gives
\(K_n\equiv^k_{\infty\omega}K_{n'}\). A first-order sentence of rank at most
\(k\) is equivalent to an \(\mathrm{FO}^k\)-sentence
(\cref{thm:p03-w03-p03-055-rank-variables}, as \(k\ge1\)), which is an
\(\mathcal L^k_{\infty\omega}\)-sentence; so \(K_n\equiv_kK_{n'}\).

(2) The sentence \(\sigma_{\ge k+1}\) belongs to \(\mathrm{FO}^{k+1}\) and
has rank \(k+1\); it holds in \(K_{k+1}\) and fails in \(K_k\). If
Duplicator won the \((k+1)\)-pebble game on \(K_k\) and \(K_{k+1}\),
\cref{thm:p03-w03-p03-055-pebble} would give
\(K_k\equiv^{k+1}_{\infty\omega}K_{k+1}\); if Duplicator won the
\((k+1)\)-round game, \cref{thm:p03-w03-p03-038-pointed-ef} would give
\(K_k\equiv_{k+1}K_{k+1}\). Either contradicts the behaviour of
\(\sigma_{\ge k+1}\).

(3) A graph satisfies \(\sigma_{\ge k+1}\) iff it has at least \(k+1\)
vertices. An \(\mathcal L^k_{\infty\omega}\)-sentence defining this class on
finite cliques would be false in \(K_k\) and true in \(K_{k+1}\),
contradicting \(K_k\equiv^k_{\infty\omega}K_{k+1}\) from (1). A sentence
defining the class on finite graphs defines it on finite cliques.

(4) A sentence of rank at most \(k\) defining the class on finite cliques
would separate \(K_k\) from \(K_{k+1}\), contradicting
\(K_k\equiv_kK_{k+1}\) from (1).
\end{proof}
\end{annexitem}

\begin{annexitem}{annex:p03-w03-p03-055-orders}{thm:p03-w03-p03-055-orders}{Two variables on finite linear orders}
\begin{proof}
(1) For \(j\ge1\) let \(\gamma_j\) stand for either \(\alpha_j\), with
\(z=x_1\) and \(z'=x_2\), or \(\beta_j\), with \(z=x_2\) and \(z'=x_1\);
let \(\gamma'_j\) be the other one of the pair. The definitions read
\(\gamma_1(z)\defeq z=z\) and
\(\gamma_{j+1}(z)\defeq\exists z'(z'<z\wedge\gamma'_j(z'))\), and the free
variable of \(\gamma'_j\) is \(z'\). We show by induction on \(j\), for both
choices at once, that in any linear order \(\gamma_j\) holds of \(c\) iff at
least \(j-1\) elements are below \(c\). For \(j=1\) both sides always hold.
For \(j+1\): \(\gamma_{j+1}\) holds of \(c\) iff some \(d<c\) satisfies
\(\gamma'_j\), iff (induction hypothesis) some \(d<c\) has at least \(j-1\)
elements below it. If so, those elements and \(d\) are \(j\) elements below
\(c\). If at least \(j\) elements lie below \(c\), choose \(j\) of them and
let \(d\) be the largest of these; the other \(j-1\) lie below \(d\), and
\(d<c\). The formulas use only \(x_1,x_2\), and
\(\operatorname{qr}(\gamma_j)=j-1\) by induction, so
\(\exists x_1\,\alpha_n(x_1)\) has rank \(n\). It holds iff some element has
at least \(n-1\) elements below it, iff there are at least \(n\) elements:
given \(n\) elements, the largest has the other \(n-1\) below it.

(2) For \(m<n\) the \(\mathrm{FO}^2\)-sentence \(\exists x_1\,\alpha_n(x_1)\)
holds in \(L_n\) and fails in \(L_m\). By
\cref{thm:p03-w03-p03-055-pebble} a winning family in the \(2\)-pebble
game would give \(L_m\equiv^2_{\infty\omega}L_n\); so Spoiler wins.

(3) Let \(\eta_n\defeq\exists x_1\,\alpha_n(x_1)\wedge
\neg\exists x_1\,\alpha_{n+1}(x_1)\), an \(\mathrm{FO}^2\)-sentence that by
(1) holds in a linear order iff it has exactly \(n\) elements, and put
\(\theta_S\defeq\bigvee_{n\in S}\eta_n\), a sentence of
\(\mathcal L^2_{\infty\omega}\) (the empty disjunction being \(\bottom\)). A
finite linear order satisfies \(\theta_S\) iff its size lies in \(S\). With
\(S\) the set of even positive integers, \(\theta_S\) defines even
cardinality on finite linear orders. That even cardinality is not
first-order definable on finite linear orders is the content of
\cref{def:p03-ef-parity}.
\end{proof}
\end{annexitem}

\begin{annexitem}{annex:p03-w03-p03-055-connectivity}{thm:p03-w03-p03-055-connectivity}{Connectivity and 2-colorability}
\begin{proof}
(1) The undefinability of connectivity is
\cref{thm:p03-w03-p03-038-connectivity}, which shows that Duplicator wins
the \(k\)-round game on \(C_{2m}\) and \(C_m\sqcup C_m\) for every
\(k\in\N\) and every \(m\ge(k+1)2^{k+2}\).

For 2-colorability and \(k\in\N\) choose \(m\) odd with
\(m\ge(k+1)2^{k+2}\).
\(C_{2m}\) is 2-colorable: color \(i\in\{0,\ldots,2m-1\}\) by its parity;
the edges \(\{i,i+1\}\) for \(i<2m-1\) and \(\{2m-1,0\}\) join vertices of
different parity. \(C_m\) is not 2-colorable: in a proper coloring with
colors \(0,1\), adjacent vertices differ, so along the path
\(0,1,\ldots,m-1\) the color of \(i\) is the color of \(0\) plus \(i\)
modulo \(2\); since \(m-1\) is even, \(m-1\) and \(0\) receive the same
color, but they are adjacent. So \(C_m\sqcup C_m\) is not 2-colorable.
For every \(k\) we thus have a 2-colorable finite graph and a finite graph
that is not 2-colorable on which Duplicator wins the \(k\)-round game; by
\cref{thm:p03-w03-p03-055-rank-definability}(2), with \(\mathcal C\) the
class of finite graphs, 2-colorability is not first-order definable on
finite graphs.

(2) By induction on \(n\): under \(x_1\mapsto u\), \(x_2\mapsto v\),
\(\pi_n\) holds iff there is a walk of length exactly \(n\) from \(u\) to
\(v\). For \(n=1\) this is adjacency. For \(n+1\): \(\pi_{n+1}\) holds iff
there is \(w\) adjacent to \(u\) such that, under
\(x_1\mapsto u,x_2\mapsto v,x_3\mapsto w\), some \(u'\) with \(u'=w\)
satisfies \(\pi_n\) under \(x_1\mapsto u'\), \(x_2\mapsto v\); that is, iff
some neighbor \(w\) of \(u\) has a walk of length \(n\) to \(v\), iff
there is a walk of length \(n+1\) from \(u\) to \(v\). Every \(\pi_n\) uses
only \(x_1,x_2,x_3\) and has free variables \(x_1,x_2\). A graph is
connected iff any two vertices are equal or joined by a walk of some length
\(n\ge1\), which is what the displayed sentence of
\(\mathcal L^3_{\infty\omega}\) says.

(3) \(C_{2m}\) is connected, since from any vertex \(i\) the walk
\(i,i+1,i+2,\ldots\) reaches every vertex, and \(C_m\sqcup C_m\) is not,
since no edge joins the two copies. So the sentence of (2) holds in
\(C_{2m}\) and fails in \(C_m\sqcup C_m\).
If Duplicator won the \(3\)-pebble game on them,
\cref{thm:p03-w03-p03-055-pebble} would make them
\(\equiv^3_{\infty\omega}\)-equivalent; so Spoiler wins.
\end{proof}
\end{annexitem}

\begin{annexitem}{annex:p03-w03-p03-055-bipartite}{thm:p03-w03-p03-055-bipartite}{2-colorability in three-variable infinitary logic}
\begin{proof}
A \emph{walk} of length \(n\ge0\) from \(u\) to \(v\) is a sequence
\(u=w_0,w_1,\ldots,w_n=v\) of vertices with \(E(w_{j-1},w_j)\) for
\(1\le j\le n\); it is \emph{closed} if \(n\ge1\) and \(u=v\). Reversing a
walk gives a walk, because \(E\) is symmetric, and two walks, the first
ending where the second starts, concatenate to a walk whose length is the
sum.

Let \(c\) be a proper 2-coloring, with colors \(0,1\). Along a walk, adjacent
vertices receive different colors, so by induction on \(j\),
\(c(w_j)\equiv c(w_0)+j\pmod2\). For a closed walk of length \(n\) this gives
\(c(w_0)\equiv c(w_0)+n\), so \(n\) is even.

Conversely suppose there is no closed walk of odd length. Write
\(u\approx v\) when a walk from \(u\) to \(v\) exists; by the remarks on
reversal and concatenation, and the walks of length \(0\), this is an
equivalence relation. Using the axiom of choice, pick a vertex \(r_Q\) in
each \(\approx\)-class \(Q\) (for a finite graph, the first vertex of \(Q\)
in a fixed enumeration). For \(v\in Q\) let \(c(v)\) be the parity of the
length of a walk from \(r_Q\) to \(v\). This does not depend on the walk: if
walks of lengths \(p\) and \(q\) lead from \(r_Q\) to \(v\), the first
followed by the reverse of the second is a walk from \(r_Q\) to \(r_Q\) of
length \(p+q\); if \(p+q\ge1\) it is closed, hence of even length, and if
\(p+q=0\) then \(p=q=0\); in both cases \(p\equiv q\pmod2\). If \(E(u,v)\),
then \(u\approx v\), and a walk of length \(p\) from \(r_Q\) to \(u\)
followed by \(v\) is a walk of length \(p+1\) to \(v\); so
\(c(v)\neq c(u)\), and \(c\) is a proper 2-coloring.

By part~(2) of \cref{thm:p03-w03-p03-055-connectivity}, under
\(x_1\mapsto u\), \(x_2\mapsto v\) the formula \(\pi_n\) (for \(n\ge1\))
holds iff there is a walk of length \(n\) from \(u\) to \(v\). So
\(\exists x_1\exists x_2(x_1=x_2\wedge\bigvee_{n\text{ odd}}\pi_n(x_1,x_2))\)
holds in a graph iff some vertex has a closed walk of odd length through
it, and its negation holds iff the graph is 2-colorable. Every \(\pi_n\)
uses only \(x_1,x_2,x_3\), so the sentence lies in
\(\mathcal L^3_{\infty\omega}\).

Let \(m\ge3\) be odd. In \(C_{2m}\), coloring \(i\in\{0,\ldots,2m-1\}\) by its
parity is proper, since every edge \(\{i,i+1\}\) with \(i<2m-1\), and the edge
\(\{2m-1,0\}\), joins vertices of different parity. In each copy of \(C_m\)
the sequence \(0,1,\ldots,m-1,0\) is a closed walk of odd length \(m\). So
the sentence above, which lies in \(\mathcal L^3_{\infty\omega}\), is true in
\(C_{2m}\) and false in \(C_m\sqcup C_m\). If Duplicator had a winning family
for the \(3\)-pebble game on them, \cref{thm:p03-w03-p03-055-pebble} would make
them \(\equiv^3_{\infty\omega}\)-equivalent; so Spoiler wins.
\end{proof}
\end{annexitem}
